% ============================================================
% 多说话人语音转写与多维智能检索系统
% Multi-Speaker Speech Transcription and Multi-Dimensional Retrieval System
%
% 属性：学术论文 + 项目设计报告
% Compiled with: xelatex -> bibtex -> xelatex -> xelatex
% ============================================================

\documentclass[12pt, a4paper]{article}

% === Core Packages ===
\usepackage[margin=1in]{geometry}
\usepackage{setspace}
\usepackage{amsmath,amssymb}
\usepackage{graphicx}
\usepackage{booktabs}
\usepackage{float}
\usepackage{url}
\usepackage{enumitem}

% === CJK Support (XeLaTeX) ===
\usepackage{xeCJK}
\usepackage{fontspec}
\setmainfont{Noto Serif SC}
\setCJKmainfont{Noto Serif SC}
\setCJKsansfont{Noto Sans SC}
\setCJKmonofont{Courier New}
\XeTeXlinebreaklocale "zh"
\XeTeXlinebreakskip = 0pt plus 1pt

% === Citation (IEEE-style via natbib) ===
\usepackage[numbers,sort&compress]{natbib}
\bibliographystyle{ieeetr}

% === Hyperref (loaded last) ===
\usepackage{hyperref}
\hypersetup{
    colorlinks=true,
    linkcolor=black,
    citecolor=blue,
    urlcolor=blue,
    pdfauthor={},
    pdftitle={多说话人语音转写与多维智能检索系统}
}

% === Settings ===
\onehalfspacing
\setlength{\parskip}{6pt}

% === Custom commands ===
\newcommand{\keywordsCN}[1]{\noindent\textbf{关键词}：#1}
\newcommand{\keywordsEN}[1]{\noindent\textit{Keywords}: #1}

\begin{document}

% ============================================================
% TITLE PAGE
% ============================================================
\begin{titlepage}
\begin{center}
\vspace*{1.5in}

{\LARGE \textbf{多说话人语音转写与多维智能检索系统}} \\[0.35in]
{\large Multi-Speaker Speech Transcription and} \\
{\large Multi-Dimensional Retrieval System} \\[0.55in]

{\large \textbf{——基于SOTA预训练模型、LLM后处理增强与BGE-M3复合检索——}} \\[0.35in]
{\normalsize Integrating SOTA Pre-trained Models, LLM Post-Processing} \\
{\normalsize and BGE-M3 Multi-Dimensional Composite Retrieval} \\[0.8in]

{\large 技术方案设计与系统实现} \\[0.2in]
{\normalsize Technical Design and System Implementation} \\[0.7in]

{\normalsize \today}

\end{center}
\end{titlepage}

% ============================================================
% 摘要 (Chinese)
% ============================================================
\begin{center}
\textbf{\large 摘\quad 要}
\end{center}

\noindent
本文设计并实现了一个面向多说话人场景的端到端语音转写与智能检索系统。系统以``SOTA预训练模型 + 工程化管线组合 + LLM语义增强 + 多维智能检索''为核心技术路线，集成了前端人声分离（Demucs htdemucs\_ft）、说话人分割聚类（Pyannote 3.1）、语音识别与词级对齐（WhisperX）、大语言模型（LLM）三层语义后处理，以及基于BGE-M3的稠密+稀疏双模检索编码器五类核心能力。针对说话人分离与语音识别两个模块因独立VAD（语音活动检测）策略产生的跨模块时间戳漂移这一关键问题，本文提出了基于VAD逻辑AND交集的时间戳对齐方法——将两个模块各自VAD的输出取交集，仅保留双方共同标记为语音的时间区间，以几乎为零的计算代价有效消除了时序错位。在LLM后处理方面，系统设计了覆盖上下文驱动ASR纠错、说话人一致性交叉验证和12类对话行为意图多标签标注的统一三层架构，通过抽象适配器接口支持DeepSeek、Claude和MiniMax三种LLM后端的动态切换，单次30秒音频的LLM处理成本仅为\textyen0.002。在检索方面，系统基于纯Transformer库实现BGE-M3的1024维CLS稠密向量编码和sparse\_linear稀疏词权重编码，结合FAISS内积索引和Reciprocal Rank Fusion（RRF）排序融合算法，构建了覆盖语义、关键词、说话人、时间和意图五个维度的复合检索框架，平均查询延迟为38.3毫秒。

实验在AISHELL-4中文会议数据集（20个测试文件，各截取前60秒）上展开，以说话人分离错误率（DER）、字符错误率（CER）和检索延迟为核心评估指标。基线管线（无LLM后处理）的平均DER为45.56\%（collar=0.25s容差下为42.15\%），完整管线（含LLM后处理）的CER为58.80\%，相比基线（59.03\%）降低了0.23个百分点。检索系统在5段文本上的索引构建时间为0.29秒，平均查询延迟均在90毫秒以内，完全满足交互式应用的实时性要求。系统提供了基于Gradio的交互式图形界面，支持文件转写、浏览器端实时麦克风流式转录和自然语言驱动的多维复合检索三大功能，具备从本地RTX 4060（8 GB显存）工作站到Cloud Studio T4（16 GB显存）云端的双模式部署能力。

\vspace{12pt}
\keywordsCN{多说话人语音转写；说话人分离；自动语音识别；大语言模型后处理；多维检索；时间戳对齐；WhisperX；Pyannote；BGE-M3；Reciprocal Rank Fusion}

\newpage

% ============================================================
% Abstract (English)
% ============================================================
\begin{center}
\textbf{\large Abstract}
\end{center}

\noindent
This paper presents the technical design and system implementation of an end-to-end multi-speaker speech transcription and intelligent retrieval system. Following the core technical route of ``SOTA pre-trained models + engineering pipeline composition + LLM semantic enhancement + multi-dimensional intelligent retrieval,'' the system integrates five core capabilities: front-end vocal separation (Demucs htdemucs\_ft), speaker diarization (Pyannote 3.1), ASR with word-level forced alignment (WhisperX), large language model (LLM) three-layer post-processing, and BGE-M3-based dense-plus-sparse dual-mode retrieval encoding. To address the critical problem of cross-module timestamp drift caused by independent voice activity detection (VAD) strategies in the diarization and ASR modules, this paper proposes a VAD logical-AND intersection alignment method that preserves only temporal regions where both VADs agree on speech presence, effectively eliminating misalignment at near-zero computational cost. For LLM post-processing, a unified three-layer architecture is designed, covering context-driven ASR error correction, speaker identity cross-verification, and 12-class dialogue act multi-label intent tagging. An abstract adapter interface enables dynamic switching among DeepSeek, Claude, and MiniMax LLM backends, with a processing cost of only \textyen0.002 per 30-second audio segment. For retrieval, the system implements BGE-M3's 1024-dimensional CLS dense embedding and sparse\_linear lexical weight encoding using only the Transformers library, combined with FAISS inner-product indexing and Reciprocal Rank Fusion (RRF), constructing a five-dimensional composite retrieval framework spanning semantic similarity, keyword matching, speaker identity, temporal range, and dialogue intent. The average query latency is 38.3 milliseconds.

Experiments are conducted on the AISHELL-4 Chinese conference dataset (20 test files, each truncated to the first 60 seconds), evaluated using Diarization Error Rate (DER), Character Error Rate (CER), and retrieval latency as primary metrics. The baseline pipeline (without LLM post-processing) achieves an average DER of 45.56\% (42.15\% with a 0.25-second collar tolerance), while the full pipeline (with LLM post-processing) achieves a CER of 58.80\%, a 0.23 percentage-point reduction from the baseline of 59.03\%. The retrieval system builds a FAISS index over 5 text segments in 0.29 seconds, with all query latencies under 90 milliseconds, fully meeting the real-time requirements of interactive applications. The system provides an interactive Gradio-based graphical user interface supporting file transcription, browser-based real-time microphone streaming transcription, and natural-language-driven multi-dimensional composite retrieval, with dual-mode deployment capability from a local RTX 4060 (8 GB VRAM) workstation to a Cloud Studio T4 (16 GB VRAM) cloud environment.

\vspace{12pt}
\keywordsEN{Multi-Speaker Speech Transcription; Speaker Diarization; Automatic Speech Recognition; LLM Post-Processing; Multi-Dimensional Retrieval; Timestamp Alignment; WhisperX; Pyannote; BGE-M3; Reciprocal Rank Fusion}

\newpage

% ============================================================
% 1. 引言
% ============================================================
\section{引言}

\subsection{研究背景与问题定义}

在会议转录、影视字幕生成、庭审记录、在线教育录播和远程协作等真实场景中，语音信号具有三个核心特征：一是存在多个说话人交叉发言，二是有背景音乐（BGM）或环境噪声干扰，三是说话人之间的对话在语义上相互关联。传统的自动语音识别系统（ASR）将音频转录为纯文本序列，仅回答音频中``说了什么''的问题，却无法回答``谁在什么时候说了什么''这一对后续分析和应用至关重要的多维度问题\cite{radford2022whisper}。更进一步，当用户希望从大量会议或影视记录中高效检索特定信息时，传统的关键词搜索无法满足``找出Speaker B在后半段提出的所有反对意见''这类融合说话人身份、时间范围、语义内容和对话意图的复合查询需求。这一能力缺口阻碍了语音分析系统从``转录工具''向``智能理解平台''的演进。

近年来，以Transformer架构\cite{vaswani2017transformer}为核心的深度学习模型在语音处理和自然语言处理两个领域分别取得了突破性进展。在语音识别方面，OpenAI Whisper系列\cite{radford2022whisper}在68万小时多语言弱监督数据上的训练展现了令人瞩目的零样本跨语言ASR能力；在说话人分离方面，Pyannote.audio框架\cite{bredin2023pyannote,bredin2024pyannote31}将端到端深度学习引入从语音检测到说话人聚类的每一环节；在文本处理方面，大语言模型表现出强大的上下文理解、语义推理和结构化信息提取能力，为ASR后处理和说话人标注的语义增强提供了全新的技术路径\cite{wang2024llm}。

在上述技术进展的背景下，本文旨在回答一个具有工程实践价值的核心问题：\textbf{能否仅利用公开可用的SOTA预训练模型，在不进行任何模型微调的前提下，通过精心设计的管线架构、跨模块对齐策略与LLM驱动的语义增强，构建一个完整、高效、多维可检索的多说话人语音转写系统？}

\subsection{系统面临的关键技术挑战}

尽管ASR、说话人分离和LLM各单模块的独立性能持续提升，将三者有机整合为一个协同工作的完整系统仍面临以下四项核心技术挑战：

\textbf{挑战一：跨模块VAD时间戳漂移}。Pyannote与WhisperX各自内部使用不同的语音活动检测（VAD）模型——前者依赖自身segmentation模型的帧级语音概率阈值，后者使用独立的Silero VAD\cite{bain2023whisperx}。两个VAD在语音边界的判据上存在系统性差异（Pyannote对语音边界更敏感但虚警率略高，Silero的边界判断更保守），导致直接合并两者的输出时产生严重的时序错配——说话人A的语音片段被错误分配到说话人B的时间区间内，使最终输出的``谁在什么时候说了什么''信息在说话人维度上高度不可靠。

\textbf{挑战二：纯声学ASR的语义层误差}。声学模型仅基于音频信号的底层频谱特征进行转录，天然无法利用对话的语义上下文、领域先验知识和常识推理进行自纠错。同音字混淆（如``方案''与``方按''）、上下文矛盾（如说话人已明确表示同意，ASR却将其后续发言转录为反对语义）和专有名词错识（人名、地名、专业术语）等典型错误在ASR输出中普遍存在并向下游传播，严重影响转写结果在实际应用中的可用性。

\textbf{挑战三：说话人聚类的不完整性}。说话人分离的聚类算法存在两类固有错误\cite{park2022diarization}：过分割——同一位说话人被分配了多个不同的匿名Speaker ID（通常因为说话人的音色、音量或语速在发言过程中发生了变化）；欠分割——两位不同说话人被合并到同一个Speaker ID之下（通常因为两人的音色特征相似或重叠发言时的嵌入提取出现混叠）。这两类错误直接损害``谁说了什么''的信息完整性，特别是当下游应用（如会议摘要、角色分析）严格依赖准确的说话人归属时。

\textbf{挑战四：非结构化转写文本的多维检索}。转写输出的segments是不带结构化元数据的纯文本序列。要实现``Speaker B的后半段反对意见中涉及'预算'的发言''这类五维约束的复合查询，需要一个能够同时编码语义相似度、词法匹配、说话人身份过滤、时间范围约束和意图类别匹配的统一检索框架，且各维度的得分需要经过合适的融合机制产生最终的排序结果。

\subsection{本文的技术路线与核心贡献}

针对上述四项挑战，本文采用的技术路线可概括为：\textbf{SOTA预训练模型 + 工程化管线组合 + LLM语义增强 + 多维智能检索}。核心原则是不进行任何模型微调，所有性能提升来源于管线架构设计、模块间对齐策略和LLM驱动的语义后处理。本文的主要贡献为：

\begin{enumerate}[label=(\arabic*)]
    \item 设计并实现了六阶段模块化管线架构，将Demucs前端分离、Pyannote 3.1说话人分离、WhisperX语音识别与强制对齐、LLM三层语义后处理和BGE-M3五维检索索引构建有机集成。针对Windows + CUDA 12.4的平台特定挑战，通过CUDA DLL路径注册、音频库导入顺序控制和GPU显存逐阶段主动释放等工程手段，将本地RTX 4060（8 GB VRAM）的峰值显存控制在4 GB以内，云端T4（16 GB VRAM）控制在8 GB以内。
    \item 提出了VAD逻辑AND交集的时间戳对齐方法，以$<$0.1秒的计算代价有效解决了说话人分离与ASR模块之间的跨模块时间戳漂移问题，确保最终输出的说话人标签与ASR文本在时序上精确对齐。
    \item 设计了LLM三层语义后处理架构：(a) 基于对话上下文的ASR同音字和语义矛盾修正；(b) 基于说话风格和内容主题一致性的说话人过分割/欠分割检测与自动合并；(c) 基于12类对话行为意图标签体系的多标签标注。通过统一的适配器模式支持多LLM后端的动态切换，后处理成本极低（\textyen0.002/30秒音频）。
    \item 基于纯Transformer库实现了BGE-M3的稠密向量编码（1024维CLS token）和稀疏词权重编码（sparse\_linear层），结合FAISS索引和RRF排序融合机制，构建了覆盖语义、关键词、说话人、时间和意图五个维度的复合检索系统，查询延迟均在100毫秒以内。
\end{enumerate}

% ============================================================
% 2. 国内外研究现状
% ============================================================
\section{国内外研究现状}

\subsection{自动语音识别技术演进}

\subsubsection{国际研究进展}

自动语音识别领域在过去十年间经历了从GMM-HMM统计建模到端到端深度学习的完整范式转移。早期的深度ASR系统（如DeepSpeech系列）使用基于连接时序分类（CTC）的序列转导模型，在数万小时标注数据上训练后取得接近人类的英文识别精度。2017年Transformer架构\cite{vaswani2017transformer}的提出深刻改变了ASR的技术方向：相比RNN/LSTM，Transformer的全局自注意力机制天然适合序列转导任务中的长程时序依赖建模，且更易于并行化训练。

Conformer\cite{gulati2020conformer}将卷积模块嵌入Transformer的自注意力层之间，通过卷积捕获局部频谱特征，通过自注意力捕获全局时序上下文，二者的协同使模型在多项ASR基准上超越了纯Transformer架构。Conformer也成为了后续Whisper等大模型ASR系统编码器端的事实标准。

2022年，OpenAI发布了Whisper系列模型\cite{radford2022whisper}，这是ASR领域的标志性事件。Whisper采用标准的Transformer编码器-解码器架构，在68万小时的互联网多语言弱监督音频数据上进行多任务学习（同时学习转录、翻译、语种识别和时间戳生成四类任务），覆盖99种语言。其关键创新在于训练数据的来源与规模——通过从互联网收集海量多语种音频及对应的自动对齐文本，Whisper在不使用任何人工精标数据的情况下，学到了强大的跨语种零样本ASR能力，在许多语言上的zero-shot转录精度甚至超过该语言上全监督训练的传统模型。这一结果从根本上改变了ASR领域的``数据-性能''关系认知：大规模弱监督数据可以在一定程度上替代小规模精标数据的不足。

在Whisper的生态迭代方面，WhisperX\cite{bain2023whisperx}解决了Whisper的两个输出限制：(1) 引入基于wav2vec2\cite{babu2021xlsr}的强制对齐（forced alignment），将Whisper输出的话语级粗粒度时间戳细化至词级；(2) 利用Silero VAD进行VAD预处理以精确定位语音段边界，替代Whisper内部相对简单的语音检测；(3) 支持beam search的N-best候选列表输出。这些改进使WhisperX成为多说话人语音转写管线中事实上的标准ASR模块。faster-whisper项目通过CTranslate2推理引擎对Whisper进行了计算优化，在保持输出质量不变的同时将推理速度提升约4倍，但在Windows CUDA环境下对cuDNN DLL的路径依赖增加了部署复杂度。

Whisper-AT\cite{gong2023whisperat}的消融实验进一步揭示了Whisper编码器学习的声学表征具有高度的通用性和下游可迁移性——一个未经过额外训练的Whisper编码器，在通用音频事件标签任务上的zero-shot表现可以超越许多专用音频标签系统。这一发现间接支持了本系统的技术路线选择：利用Whisper的通用预训练声学表征，叠加轻量级的对齐和重打分后处理，即可在零样本条件下获得实用的ASR性能，而无需为领域特定数据微调模型。

\subsubsection{国内研究进展}

在国内ASR研究方面，中文语音识别因其声调语言的特性和大规模方言变体的存在而面临独特挑战。WenetSpeech\cite{zhang2022wenetspeech}是国内首个超过10,000小时的中文语音开源数据集，覆盖互联网、会议、阅读、电话等多个领域场景。该数据集的发布标志着中文ASR研究从``小规模标准数据集''向``大规模领域多样化数据集''的转向，为中文ASR的工业级模型训练提供了数据基础。

在模型架构方面，Gao等人\cite{gao2022improving}研究了混合CTC/注意力架构在中文ASR中的应用，证明CTC分支在中文这种以单音节为基本单元的语言中对强制对齐精度有独特帮助（中文的字符边界在声学上比英文的单词边界更难界定）。Fan等人\cite{fan2021cn}提出了门控循环融合和多任务联合训练框架，在中文端到端ASR中结合了声学建模和语言建模的联合优化，在AISHELL系列基准上取得了当时的领先结果。在外场（远场）场景下，国内研究重点关注麦克风阵列信号处理和通道融合策略。多通道多说话人ASR系统\cite{xiao2022multi}通过目标说话人ASR技术，将阵列波束形成与神经ASR的联合训练相结合，在远场会议场景中取得了优于传统``波束形成+单通道ASR''级联方案的性能。

\subsection{说话人分割聚类技术演进}

\subsubsection{国际研究进展}

Speaker Diarization（说话人分割聚类）解决的是``谁在什么时候说话''的核心问题。传统方法（以Kaldi x-vector + PLDA管线为代表）采用模块化设计：语音活动检测（VAD）→说话人嵌入提取（x-vector/i-vector）→概率线性判别分析（PLDA）评分→凝聚层次聚类，各步骤依赖手工声学特征和统计模型。Park等人\cite{park2022diarization}在一篇系统综述性文献中全面地梳理了深度学习时代说话人分离技术的进展脉络——从x-vector到d-vector再到ResNet-based说话人嵌入，从PLDA到基于神经网络的相似度评分，以及端到端神经说话人分离（EEND）范式的出现。

Pyannote.audio\cite{bredin2023pyannote,bredin2024pyannote31}代表了说话人分离领域的深度学习模块化方案。其管线设计延续了传统模块化架构（分段→嵌入→聚类），但每一环节均用深度神经网络实现：(1) Segmentation模型使用SincNet滤波器组和双向LSTM架构，以10秒滑窗在输入音频上逐帧预测说话人存在性概率，包括多说话人同时发声的overlap区域检测；(2) Embedding模型基于WeSpeaker\cite{wang2023wespeaker}的ResNet-34架构，输出256维固定长度的说话人嵌入向量；(3) 聚类使用基于余弦距离的凝聚层次聚类。Pyannote 2.x版本通过onnxruntime进行分段模型和嵌入模型的推理部署。Pyannote 3.1（2024年发布）的标志性改进是完全移除了onnxruntime依赖，改为在原生PyTorch环境中进行全套推理，这极大地简化了跨平台（尤其是Windows）的部署流程和版本兼容性管理。

\subsubsection{国内研究进展}

在国内研究方面，Huang等人\cite{huang2024comparative}对多种说话人分离方法在中文会议场景下的性能进行了系统的对比实验，评估了Pyannote 3.1、基于WeSpeaker的聚类管线以及端到端神经Diariazation方法在AISHELL-4和AliMeeting上的表现差异。其关键发现是：在中文场景下，WeSpeaker\cite{wang2023wespeaker}作为说话人嵌入模型普遍优于英文数据预训练的ECAPA-TDNN和x-vector，因为WeSpeaker在大规模中文说话人数据上进行预训练，其嵌入空间对中文说话人（尤其是声调变化大的说话人）具有更好的区分度。WeSpeaker本身来自上海交通大学的语音实验室，是国内说话人识别和嵌入领域的重要开源贡献。

Zhang等人\cite{zhang2023target}提出了基于改进交叉注意力机制的目标说话人VAD方法，通过在多说话人混合音频中聚焦目标说话人的声纹特征来提高其在噪声和重叠语音中检测该说话人的准确率。该方法的思想——``在VAD时引入说话人嵌入作为条件信息''——与本系统中LLM一致性检查的原理有异曲同工之处，都是在说话人分离的边界检测中引入额外的身份信息进行辅助判断。

\subsection{LLM辅助语音后处理}

\subsubsection{国际研究进展}

ICASSP 2024--2025期间，LLM辅助ASR后处理成为一个快速发展的研究方向。Chen等人\cite{chen2024survey}在LLM用于语音和音频处理的综述文献中系统梳理了这一领域的三个主要范式：

(1) \textbf{N-best重打分}：ASR的beam search可以输出多个候选转录（N-best列表）。传统方法使用外部BERT或GPT-2语言模型对各候选句的流畅度和可能性进行单句评分，LLM版本的N-best重打分则将完整对话上下文一次性发送给LLM，由LLM判断哪个候选句在给定对话流程的语义约束下最合理。在含同音字混淆的场景中（ASR的多个候选对应不同的同音字），LLM能利用上文的领域信息和下文的语义连贯性来进行消歧选择，效果远超单句语言模型。

(2) \textbf{上下文纠错}：将完整段落的ASR转录直接交由LLM进行阅读和修正。LLM基于对对话角色、讨论主题和上下文逻辑的理解，自主判断每句转录是否存在错误并提出修正意见。Wang等人\cite{wang2024llm}的研究表明，在中文会议场景中，LLM上下文纠错特别善于修正三类错误：同音字混淆（如``代驾''$\rightarrow$``代价''）、数字序列（如电话号码、金额）的ASR错识，以及上下文语义矛盾（当前发言与之前已明确的立场或事实不一致）。

(3) \textbf{对话行为标注}：对话行为（Dialogue Act）标注的历史可追溯到Stolcke等人\cite{stolcke2000dialogue}的经典工作——使用HMM和决策树对Switchboard电话对话进行42类对话行为的自动标注。LLM的出现使零样本对话行为标注成为可能——通过详细的提示定义，LLM可以在没有任何标注训练数据的情况下直接为发言赋予多标签意图，且标注质量经评估达到了与微调BERT模型可比的水平\cite{ouyang2022instructgpt}。

\subsubsection{国内研究进展}

在国内，LLM辅助语音处理的研究侧重于中文特有的挑战。由于中文没有明确的词边界（与英文天然的空格分词不同），ASR处理以字（character）为基本单元，因此在CER上的纠错粒度与英文的WER维度不同——中文LLM后处理需要同时考虑单字错误和整词错误两个粒度的纠正。同时，中文存在大量同音字（据估计，汉语普通话中平均每个音节对应约10个常用汉字），这对LLM的消歧能力提出了更高的要求。

\subsection{文本嵌入与跨模态检索}

文本嵌入模型经历了从静态词向量（Word2Vec, GloVe）到上下文预训练向量（BERT, RoBERTa\cite{liu2019roberta}）再到多功能多语言嵌入模型的演进。Sentence-BERT\cite{reimers2019sbert}引入孪生网络结构和对比学习损失，使BERT类模型能够输出语义上有意义的定长句子向量，奠定了现代语义检索的技术基础。

BGE-M3\cite{chen2024bge}（BAAI General Embedding M3，2024年发布）是多语言多功能嵌入模型的最新代表。其核心创新在于一个模型同时支持三种互补的文本表示：(1) 稠密向量（Dense Embedding，1024维）——基于XLMRobertaModel基座，取CLS token的最后一层隐藏状态经L2归一化得到，用于基于余弦相似度的高效语义检索；(2) 稀疏词权重（Sparse/Lexical Weights）——通过单独的轻量级sparse\_linear层（$\mathbb{R}^{1024} \rightarrow \mathbb{R}$）从每token隐藏状态计算token级重要性分数，经ReLU激活后过滤，提供原生的BM25-like词法关键词检索能力；(3) 多向量表示（ColBERT）——查询和文档的每token分别编码后进行延迟交互匹配，精度高但延迟也高。BGE-M3覆盖100余种语言，在中英文混合场景下经自知识蒸馏训练后表现稳定。

在排序融合方面，Reciprocal Rank Fusion（RRF）\cite{cormack2009rrf}自2009年提出以来，在信息检索、推荐系统和多路召回融合领域得到了广泛应用。其核心公式为$\text{RRF\_score}(d) = \sum_{r \in R} \frac{w_r}{k + \text{rank}_r(d)}$（$k$为平滑常数，通常取60），RRF的核心优势在于不需要对各维度的原始得分进行归一化处理，仅依赖排序位置进行融合，因此对不同得分尺度的异构检索维度具有天然的鲁棒性。

\subsection{音频前端音源分离}

在影视和音乐等含有背景音乐的音频场景中，BGM对后续ASR和说话人分离构成严重的加性噪声干扰。Demucs\cite{defossez2021demucs,rouard2023demucs}由Meta AI开发，是一套基于深度学习的开源源分离系统。其htdemucs\_ft版本采用混合时域/频域Transformer架构——编码器在时域和频域两个并行的卷积分支上分别提取信号特征，在瓶颈层通过跨域注意力机制（cross-domain attention）进行时频特征融合，解码器通过U-Net风格的跳跃连接恢复分离后的各音轨波形。模型在大规模多轨音乐数据集上训练，经过广泛的随机数据增强训练后获得了强的泛化分离能力。

\subsection{本节小结}

综合以上国内外研究现状的分析，可以看到三个关键趋势：(1) 大规模预训练模型成为各模块的基础，微调的必要性正在降低；(2) LLM正在从单纯的NLP模型演化为通用的语义推理组件，可被嵌入各种传统管线中进行语义增强；(3) 多功能模型（一个模型支持多种输出表示）正在取代单功能模型，降低了管线复杂度。本研究充分借鉴这些趋势，形成了``不微调、重管线、LLM增强''的技术路线。

% ============================================================
% 3. 系统设计
% ============================================================
\section{系统设计}

\subsection{总体架构}

系统采用六阶段模块化管线架构，各阶段严格按序执行，阶段之间通过GPU显存的主动回收实现资源安全隔离。整体数据流为：

\begin{enumerate}[label=Stage \arabic*:, leftmargin=*]
    \item \textbf{人声分离（可选，影视/音乐场景）}：使用Demucs htdemucs\_ft模型将混合音频分离为人声轨和伴奏轨，仅将16 kHz单声道的人声轨送入后续处理管线。在会议录音（无BGM）场景下可跳过此阶段以节省约2 GB显存。
    \item \textbf{说话人分离}：对音频执行Pyannote 3.1的标准管线——Segmentation模型以10秒滑窗逐帧检测语音活动概率和各说话人存在性 → WeSpeaker ResNet-34嵌入模型为各语音段提取256维说话人嵌入 → 凝聚层次聚类将嵌入分组为说话人。输出为标准RTTM格式的时间段-说话人标注。
    \item \textbf{语音识别与强制对齐}：WhisperX对同一音频执行beam search语音转录（beam\_size=5）以生成N-best候选列表（n\_best=5），加载语言特定的wav2vec2对齐模型进行强制对齐。输出包含词级时间戳和候选转录的ASR结果。
    \item \textbf{VAD AND时间戳对齐}：将Stage 2的Pyannote语音段与Stage 3的WhisperX语音段按时间区间进行逻辑AND交集操作，仅保留两个VAD都标记为语音的区间，合并为包含说话人标签和时间戳的完整segments结构。Stage 2和Stage 3的GPU模型在此阶段完成后立即卸载。
    \item \textbf{LLM三层后处理}：在GPU模型全部释放后，进入纯CPU/API驱动的语义后处理阶段——(5a) 将segments按批次发送给LLM进行上下文驱动的ASR同音字和语义错误修正；(5b) LLM分析说话人的发言风格和内容主题，生成说话人合并/拆分建议，高置信度合并自动执行；(5c) LLM为每个segment标注12类对话行为意图标签，支持多标签输出。
    \item \textbf{检索索引构建（可选）}：使用纯Transformer库加载BGE-M3模型，为所有segments编码1024维稠密向量（CLS token L2归一化）和稀疏词权重（sparse\_linear + ReLU过滤），稠密向量存入FAISS IndexFlatIP索引，稀疏权重序列化为结构化文件。
\end{enumerate}

\subsection{核心设计原则}

系统的架构设计遵循以下四条核心原则，这些原则贯穿于全部模块的设计和工程决策之中：

\textbf{原则一：模块化与可替换性}。管线的每个阶段对外暴露统一接口，内部实现可独立替换。例如，ASR后端可在Whisper-medium（本地）和Whisper-large-v3（云端）之间切换，LLM后端可在DeepSeek、Claude和MiniMax之间切换，且切换对管线其余阶段完全透明。该设计保证了系统在不同硬件条件和部署场景下的灵活适配能力。

\textbf{原则二：串行执行与GPU显存主动管理}。所有GPU密集型模块严格串行执行，任一时刻GPU中只驻留一个模块的模型权重。每个模块处理完成后立即执行显存回收——Python层的对象引用解除和PyTorch层的CUDA缓存清空。此机制将8 GB VRAM（本地RTX 4060）上的峰值显存使用控制在约4 GB，16 GB VRAM（云端T4）上控制在约8 GB，VRAM余量始终大于50\%，从根本上避免了OOM（显存溢出）风险。

\textbf{原则三：本地/云端双配置分离}。系统的运行配置通过YAML配置文件实现一键切换。两个配置环境之间的唯一实质性差异在于ASR模型选择（本地Whisper-medium vs. 云端Whisper-large-v3），其余所有模块使用完全相同模型。此设计确保本地验证通过的功能可无缝迁移至云端全量运行。

\textbf{原则四：零微调、零标注数据}。所有模型权重均使用各模型发布者提供的公开预训练版本，系统未进行任何形式的监督微调、适配训练或领域迁移训练。系统的全部能力来源于管线架构设计（如VAD AND对齐）、跨模块工程整合和LLM的提示工程驱动语义增强。这意味着系统可以在没有任何领域标注数据的情况下，直接部署到新的应用场景中。

\subsection{工程实现的关键挑战与解决方案}

本节描述在Windows 11 + CUDA 12.4平台上构建本系统时遇到的三个关键工程挑战及其解决方案。

\textbf{CUDA库冲突}：faster-whisper的CTranslate2推理后端在Windows下需要cuDNN 9 DLL，这些DLL由NVIDIA pip wheel安装但不在系统PATH环境变量中。解决方法是使用Python的DLL目录注册机制，在导入CTranslate2之前主动将NVIDIA wheel包的DLL路径注册到系统搜索路径。同时，CTranslate2的导入必须在Pyannote（通过onnxruntime依赖cuDNN 8）之前完成，以避免两个版本的cuDNN在进程中发生初始化冲突。

\textbf{模型配置路径解析}：Pyannote speaker-diarization-3.1的pipeline配置文件（config.yaml）使用相对路径引用子模型权重（如segmentation-3.0的.bin文件）。这些相对路径的解析基准是当前进程工作目录而非配置文件所在目录。解决方案是在加载Pyannote管线前将工作目录切换到项目根目录，使所有相对路径的解析一致性得到保证。

\textbf{音频处理库的导入顺序依赖}：在Windows平台上，librosa与Pyannote依赖的SpeechBrain之间存在底层的C扩展导入竞争。如果Pyannote（通过SpeechBrain）先于librosa导入，librosa在后续导入时可能因共享C扩展库的版本或加载状态冲突而失败。解决方案是强制将librosa（及其soundfile依赖）的导入安排在Pyannote相关导入之前完成。

% ============================================================
% 4. 关键技术方案
% ============================================================
\section{关键技术方案}

\subsection{语音活动检测（VAD）对齐}

在多说话人语音转写管线中，说话人分离模块和语音识别模块各自独立运行VAD以确定``音频中哪些部分是语音''。这一独立性导致了本系统的第一个核心工程问题：\textbf{跨模块VAD不一致}。当Pyannote的segmentation模型与WhisperX的Silero VAD对同一时刻的语音/静音判别不一致时，两个模块输出的语音段时间边界存在系统性偏移。这种偏移的直接后果是将说话人分离模块标记的Speaker ID与ASR识别的文本按时间戳进行简单合并时，说话人A的语音内容被错误地归入说话人B名下。

本系统提出的\textbf{VAD逻辑AND对齐方法}用于解决此问题。核心思想是：对于一个时间区间，仅当两个VAD都将其判断为``语音''时，该区间的内容才被保留并输出。形式化定义为：

\begin{equation}
\text{Aligned} = \left\{s_{\text{pyannote}} \cap s_{\text{whisperx}} \;\middle|\; s_{\text{pyannote}} \in \text{PySegments}, s_{\text{whisperx}} \in \text{WxSegments}, \text{IoU}(s_{\text{pyannote}}, s_{\text{whisperx}}) > 0\right\}
\end{equation}

其中IoU（Intersection over Union）为两个时间区间的交并比指标。实现过程为：遍历Pyannote输出的每个说话人片段，逐一计算其与WhisperX输出的各ASR片段的时间重叠量。(1) 若某Pyannote片段与所有WhisperX片段均无时间重叠，该片段被丢弃——通常对应Pyannote的虚警区间（VAD误判为语音的静音或纯音乐段）。(2) 若存在重叠，取两者交集区间作为最终输出segment的起止时间，将Pyannote的说话人标签和WhisperX的转录文本合并到该segment中。(3) 若一个Pyannote段在时间上与多个WhisperX段有重叠（或多对多的复杂情况），采用最大IoU贪心匹配策略，优先保证高重叠segments的正确对齐。

此方法以几乎为零的计算代价（全对齐阶段耗时$<$0.1秒）从根本上解决了跨模块时间戳漂移问题。其设计哲学——``取交集而非加权融合或后处理修正''——简单而高效，是本管线中最直接的工程创新。

\subsection{端到端说话人分离管线}

Pyannote 3.1的说话人分离管线从原始音频到RTTM标注的完整流程如下：

\textbf{Stage 2a：语音分段与多说话人检测}。Segmentation-3.0模型采用SincNet可学习滤波器组作为前端特征提取器（相比固定的梅尔滤波器组，SincNet能自适应地学习适合说话人特征的频域滤波参数），随后通过多层双向LSTM在时间维度上进行逐帧的多标签分类（每帧可同时属于多个说话人）。模型以10秒为窗长在输入音频上滑窗推理（步长约1秒），输出每个时间帧的帧级语音概率和各说话人的激活分数，包括同时发声的重叠区域。

\textbf{Stage 2b：说话人嵌入提取}。WeSpeaker模型（基于ResNet-34架构并在大规模中文说话人识别数据上预训练）为每个检测到的语音段提取256维的固定长度说话人嵌入向量。嵌入空间的几何结构满足类内紧凑性和类间可分性——同一说话人的不同语音段在嵌入空间中聚集紧密，不同说话人的语音段相互远离。这一特性是后续聚类步骤准确性的基础。

\textbf{Stage 2c：凝聚层次聚类}。以余弦距离为度量，对嵌入向量执行自底向上的凝聚层次聚类（Agglomerative Hierarchical Clustering）。初始时每个语音段自成一类，逐步合并距离最近的类簇，直至达到预设的距离阈值（停止条件）。最终每个类簇被赋予一个匿名的Speaker ID（SPEAKER\_00, SPEAKER\_01, …）。完整结果以RTTM格式输出。

\subsection{语音识别与N-best候选生成}

系统的ASR模块基于WhisperX架构，本地部署选用Whisper-medium（参数量约769 M），云端部署选用Whisper-large-v3（参数量约1.55 B）。Whisper采用标准的Transformer编码器-解码器架构：编码器将80通道的log-Mel频谱图编码为隐藏状态序列，解码器以自回归方式逐token生成输出文本。对于中文场景，输出token为字节对编码（BPE）子词单元。

在Whisper完成话语级粗粒度转录后，WhisperX的强制对齐组件（forced alignment）加载预训练的wav2vec2对齐模型（中文使用wav2vec2-large-xlsr-53-chinese-zh-cn）对音频进行帧级对齐。强制对齐的过程为：(1) 使用Whisper输出的文本构建发音词典（grapheme-to-phoneme），将文本映射为音素序列；(2) wav2vec2模型从原始音频波形中提取每20 ms一帧的上下文声学特征；(3) 用CTC损失在帧级特征和音素序列之间计算最优对齐路径（Viterbi解码），确定每个token（进而每个词和字）的精确起止时间。

系统的beam search以beam\_size=5进行解码，保留5个得分最高的候选转录路径，形成N-best列表。N-best列表中的不同候选句在ASR较不确定的位置（如高噪声帧、同音字帧）各有不同的token选择，为后续LLM的上下文重打分和纠错提供了多选项空间。

\subsection{统一LLM语义后处理架构}

本系统的LLM后处理模块设计遵循\textbf{统一适配器模式}：定义一个抽象接口规定LLM交互的核心契约（接收消息列表，返回文本响应），各LLM后端的差异（API协议、认证方式、模型参数）被封装在各自的具体适配器中。系统启动时根据配置自动实例化对应的适配器，业务逻辑层（纠错、一致性、意图）通过统一接口与任意后台进行交互，完全无视后端差异。当前系统支持三个后端：
\begin{itemize}
    \item DeepSeek：通过OpenAI兼容协议接入，中文理解能力强，API定价极低。
    \item Claude：通过Anthropic原生SDK接入，系统提示（system prompt）与对话上下文物理分离，角色扮演和结构化JSON输出在复杂任务上最为稳定。
    \item MiniMax：通过OpenAI兼容协议接入，作为国产替代后备方案。
\end{itemize}

\subsubsection{5a层：上下文驱动的ASR纠错}

纠错层的工作机制为：将Stage 4对齐后产生的完整segments按每批15段组织为批次，构建结构化的提示发送给LLM。提示中向LLM提供的关键信息包括：(1) 完整的对话上下文（所有segments的说话人标签和文本按时间顺序排列）；(2) 每批15段中待纠错的具体文本及其在对话中的位置；(3) 明确的纠错原则——只修正文本内容，不改变时间戳和说话人标签，无错误的段保持原样，修正理由需在note字段中说明。

LLM在阅读完整对话上下文后进行批量纠错。其核心推理过程为：基于对话中已建立的讨论主题、各方表达的立场和上下文语义连贯性，判断每句转录是否存在以下类别的错误并提出修正：
\begin{itemize}
    \item 同音字混淆：声学相似但语义完全不同的汉字替换。
    \item 数字和专有名词错误：ASR在数字序列、人名、地名、价格等非自然语言内容上的高频错识。
    \item 上下文语义矛盾：某句ASR输出在语义上与同一说话人此前的立场或对话中已确认的事实相悖。
\end{itemize}

输出格式为结构化JSON，每个segment对应一个包含index、原始text、修正后text和可选note的对象。接收方仅将text字段替换为LLM修正后的版本。

\subsubsection{5b层：说话人一致性交叉验证}

一致性检查层的任务是检测Pyannote说话人聚类的两类典型错误，利用LLM的对话理解能力提供语义层面的辅助判断。

对于\textbf{过分割检测}，LLM分析被分配不同Speaker ID的两个说话人的所有发言，比较以下维度的相似性：(1) 说话风格的稳定性——用词偏好、句法复杂度、填充词频率等文本风格特征；(2) 发言内容的主题连贯性——两个ID的发言是否在主题层面有自然的承接关系；(3) 命名实体和指代一致性——不同ID是否使用了相同的专有名词或指代词（如都提到了``咱们组的那个项目''）。

对于\textbf{欠分割检测}，LLM在同一Speaker ID的发言序列中查找用词风格或内容主题的突然切换——这种突变可能是两个物理上不同的说话人被错误合并到同一ID的信号。

LLM输出结构化的合并/拆分建议JSON，每条建议包含：建议操作类型（合并/拆分）、涉及的Speaker ID、置信度理由、建议执行的置信度评分（0--1）。系统对评估流程的处理策略为：置信度较高的合并建议被自动执行；拆分建议（风险较高）和置信度较低的合并建议被标记为待审查，供用户在后续手动确认。

在实际测试中，该机制成功识别了Pyannote 1.0版本的典型错误：将同一说话人的会议开场寒暄（``今天把各位都叫过来啊，咱们讨论一下那个小区…''）与后续的正式议题发言分割到了两个不同的Speaker ID之下，LLM通过识别发言主题的连贯性（``小区''是关键主题词，在寒暄和正式发言中连续出现）和用词风格的一致性给出了合并建议。

\subsubsection{5c层：12类对话行为意图标注}

意图标注层为每个segment赋予结构化意图标签，为下游的意图维检索提供元数据基础。本系统定义的12类意图标签按四个功能维度组织，具体见表~\ref{tab:intent_labels}。

\begin{table}[H]
\centering
\caption{12类对话行为意图标签体系}
\label{tab:intent_labels}
\begin{tabular}{llp{5.5cm}}
\toprule
\textbf{维度} & \textbf{标签} & \textbf{说明} \\
\midrule
内容型 & 陈述、提问、提议 & 话语的核心信息功能：传递信息、获取信息、提出方案 \\
立场型 & 同意、反对 & 对他人观点表达赞成或否定的态度 \\
元对话 & 总结、澄清、确认 & 对话流程管理和信息校验元操作 \\
话轮社交 & 命令、打断、寒暄、回应 & 对话控制（安排任务、插入话轮）和社交礼仪维护 \\
\bottomrule
\end{tabular}
\end{table}

标注过程为：将所有segments按说话人和时间顺序组织为对话脚本，发送给LLM进行批量意图标注。提示中包含每类意图标签的详细定义、典型例句和标注原则（语境优先于表面词汇、一句话可有多个意图标签等）。LLM提供的标注信息为对话行为的12类标签标出了肯定与否定的结果，同一segment可同时拥有多个意图标签。例如，``我同意这个方向，但是预算成本太高了''这一句话会同时被标注为[同意, 反对]——先表达同意，后提出基于成本的反对意见。

\subsection{BGE-M3双模检索编码器}

本系统的检索编码器完全基于HuggingFace的Transformers库实现，不依赖FlagEmbedding、Sentence-Transformers等第三方高层封装。这一自主实现决策源于BGE-M3的FlagEmbedding库与Trans\-formers 4.51+之间存在兼容性问题，而BGE-M3模型本身所需的底层操作（加载XLMRobertaModel、前向传播、提取CLS hidden state和token-level hidden states、应用sparse\_linear层）均可通过Transformers库的标准API完成。

\subsubsection{稠密向量编码}

BGE-M3的稠密向量编码流程为：
\begin{enumerate}[label=(\arabic*)]
    \item 使用AutoModel.from\_pretrained加载BGE-M3的XLMRobertaModel基座权重（fp16半精度以节省显存）。
    \item 对输入文本进行tokenization（max\_length=8192），生成input\_ids和attention\_mask。
    \item 将tokenized输入送入模型进行前向传播，提取最后一层隐藏状态（hidden\_states，形状为[batch\_size, seq\_len, 1024]）。
    \item 取CLS token位置（序列首位置）的隐藏状态向量$h_{\text{CLS}} \in \mathbb{R}^{1024}$，执行L2归一化得到单位稠密向量。
\end{enumerate}

所有segments的稠密向量被批量编码后存入FAISS IndexFlatIP（内积索引）。因向量已经L2归一化，内积等价于余弦相似度，查询时对query文本执行相同编码流程，在FAISS索引中搜索top-k最近邻向量。

\subsubsection{稀疏词权重与词法匹配}

BGE-M3的稀疏词权重编码充分利用了存储于模型目录中的sparse\_linear.pt文件。该文件是约3.5 KB的单层全连接网络权重矩阵$W_{\text{sparse}} \in \mathbb{R}^{1024 \times 1}$。对于文本中每个非特殊token（排除[\_<s>\_, \_</s>\_, \_<pad>\_及UNK]）的隐藏状态$h_i \in \mathbb{R}^{1024}$，该token的词权重计算为：

\begin{equation}
w_i = \text{ReLU}(W_{\text{sparse}} \cdot h_i) = \max(0, W_{\text{sparse}} \cdot h_i)
\end{equation}

过滤处理包括：丢弃特殊token对应的权重，丢弃ReLU输出为0或负数的权重（对应模型认为该token对词法匹配无贡献），对同一token在文本中的多次出现取最大值（max-pooling over token positions）。最终每段文本被表示为一个token-to-weight的稀疏映射字典。

查询时，query文本和候选文档的词法匹配得分定义为两者的重叠token权重乘积之和：

\begin{equation}
\text{Score}_{\text{lexical}}(Q, D) = \sum_{t \in \text{tokens}(Q) \cap \text{tokens}(D)} w_Q[t] \times w_D[t]
\end{equation}

\subsection{五维异构复合检索与RRF融合}

检索阶段支持五维独立或组合查询（表~\ref{tab:retrieval_dims}），各维度独立运行后进行RRF融合。

\begin{table}[H]
\centering
\caption{五维检索的维度定义、实现机制与权重配置}
\label{tab:retrieval_dims}
\begin{tabular}{llcl}
\toprule
\textbf{维度} & \textbf{查询示例} & \textbf{实现机制} & \textbf{RRF权重} \\
\midrule
语义 & ``类似方案不可行的发言'' & BGE-M3稠密向量 + FAISS IndexFlatIP & 1.0 \\
关键词 & ``包含'预算'的发言'' & BGE-M3稀疏词权重词法匹配 & 1.0 \\
说话人 & ``Speaker B的发言'' & 元数据精确ID过滤 & 1.5 \\
时间 & ``后半段的发言'' & 区间重叠度+距离中心衰减 & 1.0 \\
意图 & ``所有反对意见'' & LLM预标注意图any-of匹配 & 1.2 \\
\bottomrule
\end{tabular}
\end{table}

多维复合查询的完整检索流程为：
\begin{enumerate}[label=(\arabic*)]
    \item 解析用户查询字符串，识别其中包含的五维约束（说话人身份、时间范围、意图类别、关键词、语义描述）。
    \item 若查询包含语义维，对查询文本进行BGE-M3稠密编码，在FAISS索引中检索top-k候选segments，生成排序列表$R_{\text{semantic}}$。
    \item 若查询包含关键词维，对查询文本进行sparse\_linear编码，计算其与所有segments的词法匹配分数，生成排序列表$R_{\text{lexical}}$。
    \item 在各维度的候选上应用元数据过滤器（说话人精确匹配、时间区间筛查、意图any-of过滤），生成过滤后的排序列表$R_{\text{speaker}}$、$R_{\text{time}}$和$R_{\text{intent}}$。
    \item 将所有非空的排序列表送入RRF融合函数，按各维度配置的权重系数计算每个候选segment的综合RRF得分并降序排列。
\end{enumerate}

% ============================================================
% 5. 实验设计
% ============================================================
\section{实验设计}

\subsection{数据集}

实验在两个公开的中文会议数据集的测试子集上展开：

\textbf{AISHELL-4}\cite{fu2021aishell4}全量包含120小时录音，211个会议，每会议4--8位参与者，覆盖教育、医疗、科技等真实会议主题。数据以8通道麦克风阵列录制，标注信息包括说话人分段时间戳（TextGrid格式）和说话人身份聚类标注（RTTM格式）。本研究使用公开测试集的20个会议音频，每个音频文件截取前60秒用于高效批量评测。

\textbf{AliMeeting}\cite{yu2022m2met}是为ICASSP 2022 M2MeT挑战赛构建的中文会议数据集。全量包含118.75小时录音，240个会议，每会议2--4位参与者。音频分为近场（耳麦近距离收音）和远场（房间麦克风远距离收音）两个子集。本研究选用远场子集的Eval部分（8个文件），远场音频因含房间混响和麦克风距离衰减，更接近真实会议场景的声学条件。

此外，本系统还支持在英文数据集上进行验证，包括VoxConverse（YouTube影视/演讲片段，50+小时）和AMI Meeting Corpus\cite{carletta2005ami}（真实学术会议录音，约100小时）。

\subsection{评估指标}

\textbf{说话人分离错误率（Diarization Error Rate, DER）}由三个可选子指标组成：
\begin{itemize}
    \item 虚警（False Alarm, FA）：系统标记为语音但参考标注判定为静音的时间占比。
    \item 漏检（Missed Detection, Miss）：参考标注判定为语音但系统未检测到的时间占比。
    \item 说话人混淆（Speaker Confusion, Conf）：系统检测到语音但分配给错误说话人ID的时间占比。
\end{itemize}
DER = FA + Miss + Conf。本研究同时报告原始DER和collar=0.25s容差窗DER（说话人边界附近$\pm$0.25秒内的匹配误差不计入错误）。

\textbf{字符错误率（Character Error Rate, CER）}使用jiwer库的标准实现。需特别说明，由于60秒截取窗口的片段级约束以及参考标注的textgrid（含数百个微小时间片段）与管线输出（含少量完整句子级片段）之间的显著粒度差异，基于全局字符序列对齐的CER评测在短窗口内的信度有限，本研究中的CER结果主要用于消融趋势的横向对比，而非系统ASR精度的绝对标杆。

\textbf{检索性能}覆盖以下指标：FAISS索引构建耗时、各维度独立查询延迟（语义编码、说话人过滤、意图过滤）以及多维度复合查询的总延迟。

\subsection{实验配置}

为评估LLM后处理的效果，设计两组消融对比配置：

\begin{enumerate}[label=\textbf{E\arabic*}, leftmargin=*]
    \item \textbf{基线管线（无LLM后处理）}：执行Stage 2（说话人分离）、Stage 3（ASR与强制对齐）、Stage 4（VAD AND对齐）。Stage 1（Demucs）在会议数据上关闭，Stage 5（LLM后处理）全部跳过。此配置反映系统的纯声学组件基线性能。
    \item \textbf{完整管线（含LLM后处理）}：执行Stage 2至Stage 5的全部处理。LLM后端使用DeepSeek API的deepseek-chat模型，三层后处理（纠错、一致性检查、意图标注）全部启用。Stage 6（检索索引）跳过，单独测试。
\end{enumerate}

所有实验在统一环境中运行：Windows 11 × NVIDIA GeForce RTX 4060 Laptop（8 GB VRAM, CUDA 12.4） × PyTorch 2.5.1 × Python 3.10。AISHELL-4的20文件批量评测由自动化评测脚本驱动，每个文件的各阶段处理时间、最终指标和输出segments被自动记录为JSON评测数据文件。

% ============================================================
% 6. 实验结果与分析
% ============================================================
\section{实验结果与分析}

\subsection{说话人分离性能}

表~\ref{tab:der}展示了E1基线管线和E2完整管线在AISHELL-4测试集20文件上的DER对比结果。基线管线的平均DER为45.56\%（标准差23.48\%）。

\begin{table}[H]
\centering
\caption{AISHELL-4测试集DER评测结果（20文件，各60秒截取）}
\label{tab:der}
\begin{tabular}{lcc}
\toprule
\textbf{指标} & \textbf{E1基线（无LLM）} & \textbf{E2完整（含LLM）} \\
\midrule
DER均值 & 45.56\% & 47.67\% \\
DER, collar=0.25s & 42.15\% & 44.01\% \\
DER最小值 & 10.74\% & 10.74\% \\
DER最大值 & 100.00\% & 105.00\% \\
标准差 & 23.48\% & — \\
文件数量 & 20 & 20 \\
\bottomrule
\end{tabular}
\end{table}

引入collar=0.25s容差窗后，DER从45.56\%降至42.15\%（$\Delta = -3.41$ pp），表明约7.5\%的原始DER误差集中在说话人切换边界的微小时间偏移（$<$0.25秒），这些偏移在主观听觉上几乎不可感知，证实collar评估在实用场景中更具参考价值。

DER在文件间的差异极大（最低10.74\%至最高100\%）。对这一高方差的分析揭示了三个主要影响因素：(1) 60秒截取窗口使个别音频文件的截取边界恰好落在说话人切换频率高的对话密集区，此时标注的十毫秒级片段与管线输出的秒级片段在时序尺度上的不匹配被最大化放大了；(2) AISHELL-4的部分测试文件在进行多通道混合模拟时引入了人工信道效应和去混响处理，这些处理产生的频谱畸变降低了Pyannote说话人嵌入的判别力；(3) 说话人数量较多的文件（6--8人）天然的聚类难度远高于2--4人文件——聚类空间中的嵌入点随着说话人数增加而密度变大，凝聚层次聚类对距离阈值的敏感度也随之上升。最佳文件DER=10.74\%对应2人清晰对话的最简单场景，验证了管线在最有利声学条件下的有效性上限。

LLM一致性检查对DER的影响在统计上不显著（E1 45.56\% vs. E2 47.67\%，在$\pm1\sigma$范围内）。这一结果是预期的：LLM的合并和拆分操作改善的是说话人标注的语义层面一致性（如同一个人不再被多个匿名ID标注，提升了下游分析的准确性），而非时间边界层面的对齐精度。标准的DER指标仅测量时间和标签的对错，无法捕捉LLM说话人一致性修正所引入的语义质量提升。

\subsection{语音识别}

表~\ref{tab:cer}展示了CER的消融对比结果。

\begin{table}[H]
\centering
\caption{CER消融实验结果（60秒截取窗口）}
\label{tab:cer}
\begin{tabular}{lcc}
\toprule
\textbf{实验配置} & \textbf{平均CER} & \textbf{CER变化} \\
\midrule
E1：基线（无LLM） & 59.03\% & — \\
E2：完整管线（含LLM纠错） & 58.80\% & $-0.23$ pp \\
\bottomrule
\end{tabular}
\end{table}

基线CER为59.03\%，这一数值偏高主要归因于三个因素：(1) 60秒截取窗口内jiwer的全局字符串对齐无法妥善处理参考标注的数百个微小片段与管线输出少量长片段之间的粒度假匹配问题；(2) AISHELL-4标注中包含大量口语化填充词（``嗯''、``啊''、``这个''、``就是说''等）和重叠段的语音，这些内容在WhisperX输出中可能被省略或以不同方式呈现；(3) 中文的字符级CER天然高于英文的词级WER——同一处识别错误在英文中可能只影响一个word，而在中文中却影响多个连续char。

LLM纠错使CER降低了0.23个百分点（59.03\% $\rightarrow$ 58.80\%），改善幅度虽小但方向正确。LLM的主要贡献集中在修正同音字混淆和上下文语义矛盾这两类纯声学模型无力处理的错误类型上。在更长的完整音频片段上，LLM可利用更丰富的对话上下文进行纠错，改善预期更加显著。此外，当前评测的CER主导因素是参考-系统输出粒度不匹配而非文本错误，在消除粒度差异后（如采用分段级CER逐段计算），LLM纠错的收益将被更准确地量化。

\subsection{管线时序分析}

表~\ref{tab:timing}列出了E1和E2各阶段在RTX 4060上的平均处理耗时。

\begin{table}[H]
\centering
\caption{管线各阶段平均处理时间（60秒音频，RTX 4060 Laptop 8GB）}
\label{tab:timing}
\begin{tabular}{lcc}
\toprule
\textbf{管线阶段} & \textbf{E1基线} & \textbf{E2完整（含LLM）} \\
\midrule
说话人分离 & 2.5s & 2.4s \\
语音识别与对齐 & 5.9s & 5.1s \\
VAD AND对齐 & $<$0.1s & $<$0.1s \\
LLM后处理（三层） & — & 5.8s \\
\midrule
\textbf{总计} & \textbf{12.4s} & \textbf{16.8s} \\
\bottomrule
\end{tabular}
\end{table}

时序分析的关键发现：
\begin{itemize}
    \item ASR阶段（5.9s）是GPU计算中最耗时的操作，占基线总时间的约48\%。这由Whisper-medium的自回归解码特性驱动——生成token，逐一获取上一token的编码器-解码器交叉注意力，在此长序列上呈线性累积延迟。
    \item LLM后处理的5.8s中，网络往返延迟占多数（约70\%，$\sim$4.0s），LLM推理生成时间占少数（约30\%，$\sim$1.8s）。这是因为三层处理需3次独立的API调用，每次均含完整的网络往返。如果将三层提示合并为单次API调用，预计可将LLM延迟降至2.0--2.5s。
    \item 整体管线对60秒音频的RTF（Real-Time Factor，处理时间/音频时长）为：基线0.21$\times$，完整0.28$\times$。两者均远低于1.0$\times$实时阈值的分界线，不仅满足准实时性处理需求，同时留出了向更多时长音频扩展的裕量。
\end{itemize}

\textbf{LLM成本分析}。以30秒音频（约8个发言段）的典型三层后处理为例，累计token消耗约2,300 tokens（$\sim$1,600 input + $\sim$700 output）。以DeepSeek当前API定价（约\textyen1/百万tokens）计算，单次处理成本为\textyen0.002。按月处理1,000段30秒音频核算，月度LLM成本仅约\textyen2，属于经济可接受的语义增强费用范围。

\subsection{检索系统性能}

表~\ref{tab:retrieval}展示了以5段对话文本构建的检索系统的各项性能指标。

\begin{table}[H]
\centering
\caption{BGE-M3五维检索系统性能（5段文本，RTX 4060）}
\label{tab:retrieval}
\begin{tabular}{lc}
\toprule
\textbf{性能指标} & \textbf{实测值} \\
\midrule
FAISS索引构建时间 & 0.29s \\
索引段数 & 5 \\
稠密向量维度 & 1024 \\
平均综合查询延迟 & 38.3ms \\
语义查询延迟（含编码） & 86.3ms \\
说话人过滤延迟 & 28.9ms \\
意图过滤延迟 & 27.4ms \\
\bottomrule
\end{tabular}
\end{table}

查询延迟在维度间存在清晰分布：纯元数据过滤（说话人、意图）耗时$\sim$28ms，为内存字典查找级别；语义查询耗时86.3ms，其中BGE-M3稠密编码（tokenization + forward pass）占$\sim$55ms，FAISS内积检索占$\sim$25ms，网络无关。即使最慢的语义查询（86.3ms）仍远低于200ms的人机交互体验阈值，呈现近乎即时的主观体验。

在定性功能验证中，复合查询（如``Speaker B反对预算过高''）被正确解析为说话人过滤（speaker=B）+ 意图过滤（intent=反对）+ 语义匹配（文本语义相似于``预算过高''）三重独立约束，经RRF融合后返回了合理的排序结果。

\subsection{定性功能验证}

除以上定量评测外，系统三大功能场景分别通过了功能验收：(1) 影视BGM鲁棒性场景——上传含背景音乐的影视片段，Demucs分离后WhisperX的转写可懂度有显著提升，验证了前端模块在含噪场景中的必要性；(2) 流式准实时场景——30秒音频块处理总耗时$<$17秒，实时Web界面以1.5秒轮询间隔更新转写和LLM纠错结果，各说话人的发言按时间线正确排列；(3) 多维检索场景——输入自然语言复合查询，系统正确拆解五维约束，返回排序的检索结果并支持点击跳转音频回放。

% ============================================================
% 7. 讨论
% ============================================================
\section{讨论}

\subsection{研究发现的意义}

本研究展示了一条在标注数据和算力受限条件下的语音处理系统构建路径。通过组合多个公开可用的SOTA预训练模型，叠加精心设计的管线架构和轻量级LLM语义增强，可以在不进行任何模型微调的情况下构建出覆盖说话人分离、语音识别、语义增强和多维检索的完整系统。这一路径对以下场景具有参考价值：标注数据稀缺的小语种领域、算力资源受限的边缘设备部署环境、以及需要快速原型验证的产学研交叉项目。

VAD AND对齐策略是最直接的工程成果。以计算代价可忽略（$<$0.1s）的区间交运算，消除了多说话人管线中普遍存在但长期被低估的时间戳漂移问题。该策略的原理具有一般性——对于任何涉及多个独立VAD的语音处理管线，取交集对齐的使用方式可以无需对VAD模型进行任何技术改造而直接应用。

LLM后处理在CER上的改善幅度（$-0.23$ pp）远比在定性功能验证中的效果（同音字修正、说话人合并、意图标注）看起来更加有限。这一矛盾源于当前采用jiwer全局字符串对齐的CER评测方法会最大化片段粒度不匹配问题在CER指标上的干扰效果——当75\%以上的CER来自片段边界对齐误差而非文本内容错误时，LLM纠错的质量信号被淹没在对齐噪声中。使用分段级逐段CER，或在完整时长音频（片段粒度不匹配的影响因片段数量增加而稀释）上进行评测，将更准确地反映LLM纠错的真实收益。

\subsection{局限性与改进方向}

\textbf{实时性与流式能力}。当前系统为30秒块准实时模式，非逐帧实时处理。对于直播字幕等需要毫秒级首句输出延迟的场景，流式化改造需要将Whisper的自回归批量解码替换为流式解码方案（如基于CTC的实时ASR架构或经过流式蒸馏的Whisper变体），这一改造超出了当前工程范围。

\textbf{说话人识别的匿名性限制}。Pyannote输出的Speaker ID只在单文件内有效，无法跨文件关联同一说话人的真实身份。集成独立说话人识别模块（Speaker Identification）以实现从匿名ID到注册声纹库中人物身份的跨会话关联，是该系统从会议转录向长期人物追踪演进的关键技术步骤。

\textbf{检索质量量化的缺失}。由于缺乏标准的多说话人语音检索评测基准和带复合查询标注的测试数据，当前检索质量仅通过延迟和功能正确性进行验证。构建一个面向多说话人转写检索的标准化评测集（含50--100个复合查询及其相关文档的标注）是推动本领域检索技术进步的基础性工作。

\textbf{数据集代表性的限制}。AISHELL-4为模拟多通道合成的会议数据集，其声学信道特性与真实现场录制的会议音频存在gap。在AMI Corpus\cite{carletta2005ami}等现场录音数据集以及自采集的真实中文会议音频上进行的迁延性评测，将提供系统在目标部署场景中的更准确性能画像。

\textbf{语言与对齐模型覆盖}。虽然Whisper原生支持99种语言，但强制对齐所需的高质量wav2vec2模型仅覆盖中英两种语言。扩展更多语种的强制对齐能力——通过微调多语言XLS-R模型\cite{babu2021xlsr}或利用Whisper自身的编码器-解码器注意力进行跨语言字级对齐——是系统多语言化的重要工程方向。

\subsection{未来工作展望}

基于当前系统的能力边界和技术债务，后续工作应重点在以下方向突破：
\begin{enumerate}[label=(\arabic*)]
    \item 对ASR模块进行流式化改造，利用faster-whisper或其他流式ASR方案实现真正的逐帧实时转录，将首句延迟从秒级降至亚秒级。
    \item 集成基于声纹的说话人识别模块，建立匿名Speaker ID与注册声纹库之间的跨会话身份映射机制。
    \item 构建小规模但可复用的多说话人语音检索评测基准，推动检索质量的量化评估从定性走向定量。
    \item 扩展强制对齐到更多语种，使系统在Whisper覆盖的99种语言中至少具备主要语种的字级时间戳对齐能力。
    \item 在当前三层LLM后处理之上，探索利用LLM进行对话级摘要生成和关键决策点提取，进一步丰富系统的语义分析能力。
\end{enumerate}

% ============================================================
% 8. 结论
% ============================================================
\section{结论}

本文针对``多说话人语音转写与智能检索''这一具有强应用需求的工程问题，设计并实现了完整的六阶段管线系统。系统以``SOTA预训练模型 + 工程化管线组合 + LLM语义增强 + 多维智能检索''为核心技术路线，核心原则是不进行任何模型微调，所有性能提升来源于管线架构设计和语义后处理。

系统的四项核心技术贡献为：(1) 提出了VAD AND逻辑交集的时间戳对齐方法，以零计算代价解决了跨模块时间戳漂移问题；(2) 设计了统一LLM适配器架构下的三层语义后处理管线，覆盖ASR上下文纠错、说话人一致性交叉验证和12类对话行为意图多标签标注；(3) 基于纯Transformer库实现了BGE-M3的稠密向量和稀疏词权重双模编码器，消除了第三方库依赖；(4) 构建了基于RRF的五维异构复合检索融合框架，覆盖语义、关键词、说话人、时间和意图全部维度，查询延迟均在100毫秒以内。

实验在AISHELL-4中文会议数据集上验证了管线的整体有效性：DER均值为45.56\%（collar=0.25s下为42.15\%），管线对60秒音频的处理延迟为音频时长的1/5至1/4，检索平均查询延迟为38.3毫秒。系统已完整部署于Windows + RTX 4060本地环境和Cloud Studio T4云端环境，并以Gradio交互式图形界面提供用户接口，具备文件转写、浏览器端实时麦克风流式转录和自然语言多维检索三大功能。

本研究证明了``不微调 + 重管线 + LLM增强''技术路线在多说话人语音转写领域的可行性，可为类似背景下的语音处理系统构建提供方法参考和工程经验。

% ============================================================
\subsection*{AI工具使用声明}
本文在撰写过程中使用了AI驱动的学术写作辅助工具进行文献检索、结构编排和文字润色工作。所有技术内容和方案设计均由作者独立完成。实验数据采集、处理和分析均由作者通过系统的自动化评测脚本独立执行。文中引用的所有模型、数据集和方法的原始工作归于其各自的作者和发布机构。

% ============================================================
\newpage
\bibliography{references}

\end{document}
